\subsection{记号}

\begin{enumerate}
\def\labelenumi{(\alph{enumi})}
\tightlist
\item
  作为一般规则，交换幺半群的律用加法记法书写。此时约定 -x 表示 x 的负元。记号 \(x + (-y)\) 缩写为 x - y，类似地
\end{enumerate}

\[
x + y - z, \quad x - y - z, \quad x - y + z - t, \quad \text { etc. } \dots
\]

分别表示

\[
x + y + (- z), \quad x + (- y) + (- z), \quad x + (- y) + z + (- t), \quad \text { etc. } \dots
\]

对于 \(n \in \mathbf{Z}\) ，记号 \(\top^n x\) 换为 \(nx\) 。公式（9）至（12）于是变为

\[
(m + n). x = m. x + n. x\tag{13}
\]

\[
0. x = 0\tag{14}
\]

\[
1. x = x\tag{15}
\]

\[
m. (n. x) = (m n). x\tag{16}
\]

其中 m 和 n 属于 N，甚至属于 Z（如果 x 有负元）。同样在后一种情形下，关系式 \((-1)\) 。x = -x 成立。我们还注意到公式

\[
n. (x + y) = n. x + n. y.\tag{17}
\]

\begin{enumerate}
\def\labelenumi{(\alph{enumi})}
\setcounter{enumi}{1}
\tightlist
\item
  设 \(\mathbf{E}\) 是以乘法记法书写的幺半群。对于 \(n \in \mathbf{Z}\) ，记号 \(\frac{n}{T} x\) 换为 \(x^n\) 。我们有关系式
\end{enumerate}

\[
\begin{array}{r l} x ^ {m + n} & = x ^ {m}. x ^ {n} \\ x ^ {0} & = 1 \\ x ^ {1} & = x \\ (x ^ {m}) ^ {n} & = x ^ {m n} \end{array}
\]

以及 \((xy)^n = x^n y^n\) （若 \(x\) 和 \(y\) 可交换）。

当 \(x\) 有逆元时，这恰好是 \(x^{-1}\) 。记号 \(\frac{1}{x}\) 也用来代替 \(x^{-1}\) 。最后，当幺半群 \(E\) 是交换的， \(\frac{x}{y}\) 或 \(x / y\) 也用来表示 \(xy^{-1}\) .

